\section{Completion of the field-strength construction}
\label{sec:main-theorem-closure}

\begin{proof}[Proof of \cref{thm:main}]
Choose the construction constants by \cref{prop:ordered-constants},
in the stated order.
In particular $\sigma>0$ in
\eqref{eq:transfer-charge-exponents}.
Take the positive physical ceiling $e_G$ in
\eqref{eq:constants-physical-charge}. Theorem~\extref{D}{CoverTuning} supplies real bare mass shifts
$m_{0,N}(e)$, independent of the volume and twist,
for $0<e\le e_G$; below its defining mesh threshold they
are set to zero. Their bound
\exteq{65}{CoverTuning} gives
$|m_{0,N}(e)|<1/8$ for every $N$.
The all-flux laws at \eqref{eq:cutoff-scales} are therefore
defined at every cutoff, even when the total lattice mass
has either sign.
The physical ceiling in \cref{thm:main} is
$e_{\mathrm{phys}}=e_G$.
Theorem~\extref{A}{PositiveProbe} proves the existence
of the probes in \eqref{eq:positive-coefficient}.
Fix any such probe and $0<e\le e_\beta$, with
$e_\beta$ supplied by \cref{prop:ordered-probe-constants}.

For this probe,
Proposition~\extref{11.1}{ReferenceResponse}
provides $e_\sharp(\beta)>0$.
The terminal partition bound
\exteq{5}{TerminalPressure} and the response estimate
\exteq{119}{ReferenceResponse} supply
both parts of \eqref{eq:transfer-terminal-input}.
The free-coefficient bound
\eqref{eq:transfer-free-input} supplies its remaining
reference input. Thus \cref{prop:transfer-UV-inputs}
applies throughout this interval, with one $C_U$
independent of the cutoff, volume, placement and limit law.
Its estimates hold whenever
$N\ge N_{\rm ref}(e)$ and $T\ge T^*(e)+1$.
These are finite thresholds for each fixed $e$.

\emph{Volume-first laws and their Euclidean components.}
The tightness estimate \eqref{eq:limit-uniform-tightness}
shows that every fixed-spacing volume-limit set
$\mathcal V_N$ is nonempty and that the collection of
these limits is tight on $H_{-4}$.
Consequently volume-first limit points exist.
Let $\mu$ be any one of them, and write
$\mu^{(N_i)}\to\mu$, where $N_i\to\infty$ and
$\mu^{(N_i)}\in\mathcal V_{N_i}$.
For each $i$, choose a volume index $T_i$ along a
defining subsequence for $\mu^{(N_i)}$ such that
\[
 T_i\ge i,\qquad
 d_{\rm w}(P_{N_i,T_i},\mu^{(N_i)})\le1/i,
\]
where $d_{\rm w}$ metrizes weak convergence on $H_{-4}$.
Then
$d_{\rm w}(P_{N_i,T_i},\mu)
\le1/i+d_{\rm w}(\mu^{(N_i)},\mu)\to0$.
This is the joint sequence in
\cref{prop:limit-nine-mirrors};
both indices diverge, and all fixed finite-family
moments converge along it.

\Cref{prop:limit-nine-mirrors} gives closedness,
translation invariance, Gaussian domination with constant
one, and invariance and reflection positivity at all
planes in the nine mirror families.
Sections~\extref{5--6}{ErgodicComponents} transfer
these properties to a measurable set of full measure
in the translation-ergodic decomposition.
The Gaussian bound holds there for all real Schwartz
tests at once, by the countable-density and Fatou
argument in \exteq{20}{ErgodicComponents}.
Section~\extref{5}{ErgodicComponents} transfers each additional mirror
by transporting the conditional laws and using
uniqueness of the decomposition. Intersecting the
nine full-measure sets gives all mirror forms in
each component; these reflections also generate
the cubic group.

\Cref{thm:rotation-covariance} now applies
separately to $\mu$ and to every component in this
full-measure set. It proves their invariance under
the full Euclidean group and reflection positivity
at every plane.
For an ergodic component, the sharp cone and
\exteq{15}{ErgodicComponents} identify the
translation-invariant projection with the constant
projection. Proposition~\extref{4.1}{ErgodicComponents}
therefore gives clustering for all $L^2$ observables
as $|a|\to\infty$.
\Cref{prop:tensor-OS} gives
\textup{(E0$'$)}--\textup{(E3)} for the parent law and
\textup{(E0$'$)}--\textup{(E4)} for these components,
with exactly the tensor action and growth bound of
\cref{thm:main}.
In particular, supported tensor evaluations belong
to the positive-side $L^2$ space even when their
supports meet the reflection plane, by
\eqref{eq:tensor-closed-halfspace}.
For clustering the relevant identity is
\[
 S_{k+l}^\xi(\Psi\otimes\tau_a^{\otimes l}\Phi)
   =\E_{\mu_\xi}[Q_\Psi\,U(a)Q_\Phi],
                 \qquad Q_\Psi,Q_\Phi\in L^2(\mu_\xi).
\]
It gives \eqref{eq:main-cluster} on all tensor Schwartz
tests, using \(A=\overline{Q_\Psi}\) in the
$L^2$ clustering formula.

\emph{Fourth cumulants in the same components.}
The prescribed parent $\mu$ is the volume-first limit of the
same all-flux laws, at the same volume-independent mass assignment,
used in \cref{prop:transfer-UV-inputs}. Their physical interpolation
and weak topology are those of Theorem~\extref{A}{VolumeTransfer}.
The preceding argument gives translation invariance of $\mu$
and Gaussian domination with constant one on a full set of
its translation-ergodic components. Thus the component hypothesis
\exteq{11}{VolumeTransfer} holds for this parent.

The reference, placement and charge hypotheses of that theorem
were verified in \cref{sec:uniform-transfer}. Its conclusion
\exteq{12}{VolumeTransfer} gives a full-measure component set on which
\[
 |\kappa_4^{\mu_\xi}(f)-e^4b(f)|\le C_fe^{4+\sigma},
 \qquad C_f=8C_U+3C_U^2.
\]
The prescribed-parent clause keeps the law $\mu$ fixed;
the coefficient subsequence used in its proof does not select
a different parent. Intersect this set with the full-measure
set already used for the tensor OS conditions.
The charge choice \eqref{eq:constants-probe-charge} gives
$C_fe^\sigma\le b(f)/2$, hence
$\kappa_4^{\mu_\xi}(f)\ge e^4b(f)/2>0$ on this intersection.
The constant is independent of the parent and component.
This proves \eqref{eq:main-cumulant}; the full-measure set
is asserted for the fixed charge, probe and parent law.
A centered Gaussian variable has zero fourth cumulant, so
these component laws are non-Gaussian.
\end{proof}
